% Shared look for every deck in this repo: Metropolis with a progress bar under each
% frame title. Mirrors sourdough-science/slides/preamble.tex. Build with Tectonic:
% `make` in this directory.

\usetheme[progressbar=frametitle,numbering=fraction,sectionpage=progressbar,block=fill]{metropolis}

% Metropolis looks Fira Sans up by system name; load it from the TeX bundle instead,
% Light for body text as Metropolis intends, with lining figures for data.
\usepackage[sfdefault,light,lf]{FiraSans}
\usepackage{booktabs,tabularx,array}
\usepackage{tikz}
\usetikzlibrary{arrows.meta,positioning,shapes.geometric,calc}
\usepackage{pgfplots}
\pgfplotsset{compat=newest}

% Metropolis's palette, with this repo's own colour names mapped onto it.
\definecolor{mDarkTeal}{HTML}{23373B}
\definecolor{mLightBrown}{HTML}{EB811B}
\colorlet{ink}{mDarkTeal}
\definecolor{panel}{HTML}{EDF1F2}
\colorlet{accent}{mLightBrown}
\definecolor{flag}{HTML}{B03A2E}
\definecolor{ok}{HTML}{2E7D4F}
\definecolor{muted}{HTML}{7A7F80}

% Top-align frames, so the \vfill in \source puts the source line at the bottom.
% A slimmer title bar and narrower margins than Metropolis's defaults leave room
% for the denser slides.
\makeatletter
\beamer@centeredfalse
\setlength{\metropolis@frametitle@padding}{1.4ex}
\setlength{\metropolis@progressinheadfoot@linewidth}{2pt}
\setlength{\metropolis@progressonsectionpage@linewidth}{1.5pt}
\setlength{\metropolis@titleseparator@linewidth}{1pt}
% Metropolis's title page is \paperheight tall, which overflows a top-aligned frame.
\setbeamertemplate{title page}{%
  \begin{minipage}[b][\textheight]{\textwidth}
    \vfill
    \usebeamertemplate*{title}
    \ifx\insertsubtitle\@empty\else\usebeamertemplate*{subtitle}\fi
    \usebeamertemplate*{title separator}
    \usebeamertemplate*{author}
    \usebeamertemplate*{institute}
    \usebeamertemplate*{date}
    \vfill
  \end{minipage}}
\makeatother
\setbeamersize{text margin left=1.2em,text margin right=1.2em}

\setbeamertemplate{frame footer}{\insertshorttitle}
\setbeamertemplate{footline}{%
  \begin{beamercolorbox}[wd=\paperwidth,ht=2.6ex,dp=1.2ex,leftskip=1.2em,rightskip=1.2em]{footline}%
    \usebeamerfont{page number in head/foot}%
    \usebeamertemplate*{frame footer}\hfill\usebeamertemplate*{frame numbering}%
  \end{beamercolorbox}}
\setbeamerfont{page number in head/foot}{size=\tiny}
\setbeamercolor{footline}{fg=muted}

% The conclusion of a slide, in one sentence.
\newcommand{\reading}[1]{%
  \par\medskip
  \colorbox{panel}{\parbox{\dimexpr\linewidth-2\fboxsep}{\small\textbf{\color{ink}What it means:} #1}}}

% Small source line, placed at the bottom of a frame.
\newcommand{\source}[1]{%
  \vfill{\color{muted}\tiny Source: #1}}

% Status tags. Filled, so they read on the dark title bar as well as on the body.
\newcommand{\flagtag}[2]{\mbox{\setlength{\fboxsep}{1.2pt}\colorbox{#1}{\color{white}\scriptsize\bfseries\,#2\,}}}
\newcommand{\donetag}{\flagtag{ok}{done}}
\newcommand{\presenttag}{\flagtag{ok}{present}}
\newcommand{\pendingtag}{\flagtag{accent}{pending}}
\newcommand{\ruledout}{\flagtag{muted}{ruled out}}
\newcommand{\phantomtag}{\flagtag{flag}{Code 45}}
\newcommand{\hyp}{\flagtag{flag}{hypothesis}}
\setbeamertemplate{frametitle continuation}[from second][(cont.)]

% Nodes for TikZ diagrams: a normal box, and a box for a device that is missing.
\tikzset{
  box/.style={draw=ink,fill=panel,rounded corners=2pt,align=center,
    font=\scriptsize,inner sep=4pt,text width=#1},
  box/.default=2.6cm,
  gone/.style={draw=flag,dashed,fill=white,rounded corners=2pt,align=center,
    font=\scriptsize,inner sep=4pt,text width=#1},
  gone/.default=2.6cm,
  arrow/.style={-{Stealth[length=5pt]},draw=ink,thick},
  gonearrow/.style={-{Stealth[length=5pt]},draw=flag,dashed,thick},
}

\newcolumntype{L}{>{\raggedright\arraybackslash}X}
\newcommand{\muted}[1]{{\color{muted}#1}}

\author{ZRRBITE-PC hardware notes}
\institute{\texttt{C:/dev/hardware-notes}}
